\begin{lemma}[Scheide Proposition 4.2, full-trunk case]
Let $(T,e_0,\varphi)$ be a balanced normal globally maximum $k$-triple in
a finite loopless multigraph $G$, and suppose every vertex of $G$ is
incident with an edge.  If $h(G)=t(G)$, then $G$ is elementary.
\end{lemma}
\begin{proof}
Let $p$ be the chosen height of the normal triple.  By
\noderef{BalancedGlobalMaximumTriple}, $p=h(G)$.  By
\noderef{GlobalMaximumTriple}, $|V(T)|=t(G)$, while injectivity of the
ordered vertices in \noderef{TreeSequence} and the definition
\noderef{TreeSequenceVertexSet} give $|V(T)|=\operatorname{length}(T)+1$.
The equality $h(G)=t(G)$ therefore gives
$p=\operatorname{length}(T)+1$.  Thus all vertices and all non-first tree
edges belong to the trunk recorded by \noderef{NormalGlobalMaximumTriple}.
Its closing path lies in the $(\alpha,\beta)$ incidence graph, so each
non-first tree edge is colored $\alpha$ or $\beta$.  In the terminology of
\noderef{TreeUsedColors}, no color other than $\alpha,\beta$ is used on
$T$.

We next show that every vertex of $V(T)$ misses a free color.  The first
edge of \noderef{TreeSequence} is $e_0$ and attaches the first two ordered
vertices, so those vertices are exactly the two ends of $e_0$.  By
\noderef{ScheideGlobalMaximumMissingCounts}, each of them misses at least
two colors, whereas every other tree vertex misses at least one.  The
normal-triple data say that $\alpha$ is missing at the first vertex and
$\beta$ at the second.  The set $V(T)$ is elementary by
\noderef{ScheideGlobalMaximumOrder}; hence $\beta$ is not missing at the
first vertex, $\alpha$ is not missing at the second, and neither trunk
color is missing at any other vertex.  Each end of $e_0$ therefore has a
second missing color outside $\{\alpha,\beta\}$, and every other vertex
has all its missing colors outside that set.  Since only
$\alpha,\beta$ are used on non-first tree edges, such a color lies in the
set \noderef{TreeFreeColors}.  This proves the assertion.

Suppose for contradiction that a defective color $\delta$ exists.  Choose
a tree vertex $u$ and a free color $\gamma$ missing at $u$, as supplied by
the preceding paragraph.  The node
\noderef{ScheideDefectiveChainPathExists} gives a full simple
$(\delta,\gamma)$ chain from $u$ to a vertex outside $V(T)$.  Applying
\noderef{ScheideDefectiveChainDepartureReturn} to this supplied chain gives
its first departure and its last return, in the exact sense of
\noderef{ScheideDepartureReturn}.  If $v$ is that last return, the same
witnesses $u,\gamma,\delta$, the same full chain, and its last-return index
satisfy \noderef{ScheideExitVertex}; thus
$v\in\noderef{ScheideExitVertices}(T,\varphi)$.  But
\noderef{ScheideExitNoFree} says that the missing-color set at an exit
vertex is disjoint from the free-color set.  This contradicts the fact
proved above that every vertex of $V(T)$, in particular $v$, misses a free
color.  Hence the defective-color set is empty.

Now \noderef{ScheideGlobalMaximumOrder} makes $V(T)$ elementary and
\noderef{ScheideGlobalMaximumBoundaryFacts} makes it strongly closed.  The
two ends of $e_0$ lie in $V(T)$ by the first attachment clause of
\noderef{TreeSequence}.  Finally, the globally maximum condition contains
the $k$-triple condition \noderef{KTriple}: $G$ is edge-critical,
$\chi'(G)=k+1$, $\Delta(G)+1\leq k$, and $\varphi$ is proper with exactly
$e_0$ uncolored.  Together with the stated no-isolated-vertex hypothesis,
the witnesses $e_0,\varphi,V(T)$ satisfy condition (iv) of
\noderef{egraph2}.  The implication (iv)$\Rightarrow$(i) in that theorem
proves that $G$ is elementary.
\end{proof}
